\documentclass[12pt,a4paper]{article}

\usepackage[utf8]{inputenc}
\usepackage[T1]{fontenc}
\usepackage[spanish]{babel}
\usepackage{geometry}
\usepackage[table,xcdraw]{xcolor}
\usepackage{titlesec}
\usepackage{fancyhdr}
\usepackage{booktabs}
\usepackage{array}
\usepackage{longtable}
\usepackage{tabularx}
\usepackage{multirow}
\usepackage{parskip}
\usepackage{enumitem}
\usepackage{hyperref}
\usepackage{microtype}
\usepackage{graphicx}

\geometry{top=2.5cm, bottom=2.5cm, left=3cm, right=2.5cm}

% ── Colores ──────────────────────────────────────────────────
\definecolor{ecogreen}{RGB}{27,94,32}
\definecolor{ecogreenlight}{RGB}{220,237,200}
\definecolor{darkgray}{RGB}{50,50,50}
\definecolor{lightgray}{RGB}{240,240,240}

% ── Secciones ────────────────────────────────────────────────
\titleformat{\section}
  {\color{ecogreen}\large\bfseries}{\thesection.}{0.5em}{}
  [\color{ecogreen}\titlerule]
\titleformat{\subsection}
  {\color{darkgray}\normalsize\bfseries}{\thesubsection}{0.5em}{}
\titleformat{\subsubsection}
  {\color{darkgray}\normalsize\itshape}{\thesubsubsection}{0.5em}{}

% ── Encabezado y pie ─────────────────────────────────────────
\pagestyle{fancy}
\fancyhf{}
\fancyhead[L]{\small\itshape\color{gray}EcoMarket --- UNAP Puno}
\fancyhead[R]{\small\itshape\color{gray}P\'agina \thepage}
\fancyfoot[C]{\small\color{gray}Universidad Nacional del Altiplano --- Puno, Per\'u 2026}
\renewcommand{\headrulewidth}{0.4pt}

% ── Tipos de columna ─────────────────────────────────────────
\newcolumntype{L}[1]{>{\raggedright\arraybackslash}p{#1}}
\newcolumntype{C}[1]{>{\centering\arraybackslash}p{#1}}

\hypersetup{colorlinks=true,linkcolor=ecogreen,urlcolor=ecogreen}
\setlength{\parskip}{6pt}
\setlength{\parindent}{0pt}

% ─────────────────────────────────────────────────────────────
\begin{document}
\thispagestyle{empty}

% ══ PORTADA ══════════════════════════════════════════════════
\begin{center}
\vspace*{0.8cm}
{\Large\bfseries UNIVERSIDAD NACIONAL DEL ALTIPLANO PUNO}\\[0.3cm]
{\large Facultad de Ingenier\'ia de Sistemas}\\[0.1cm]
{\normalsize Escuela Profesional de Ingenier\'ia de Sistemas}\\[1.2cm]
\noindent\rule{7cm}{2pt}\\[1.5cm]
{\fontsize{40}{48}\selectfont\bfseries\color{ecogreen}EcoMarket}\\[0.6cm]
{\Large Plataforma de Comercio Electr\'onico de Productos}\\[0.2cm]
{\Large Ecol\'ogicos con Auditor\'ia Qu\'imica mediante IA}\\[1.2cm]
\noindent\rule{7cm}{0.5pt}\\[1cm]
{\large\bfseries INFORME DE AVANCE DE PROYECTO}\\[0.2cm]
{\normalsize\itshape Examen Parcial --- Unidad 1}\\[2cm]

\begin{tabular}{rl}
\textbf{Curso:}          & Ingenier\'ia de Software II \\[0.3cm]
\textbf{Grupo:}          & Grupo A --- Octavo Semestre \\[0.3cm]
\textbf{Docente:}        & Consorcio de Ingenier\'ia de Software \\[0.3cm]
\textbf{Estudiantes:}    & Garate Quispe Josep Vladimir \\
                         & Quispe Huanca Arnold Diego \\[0.3cm]
\textbf{Fecha:}          & Mayo 2026 \\[0.3cm]
\textbf{URL producci\'on:} &
  \href{https://eco-market-final-software-s2w6.vercel.app}%
       {\texttt{eco-market-final-software-s2w6.vercel.app}} \\
\end{tabular}

\vfill
{\small Puno --- Per\'u, 2026}
\end{center}

\newpage
\tableofcontents
\newpage

% ══ 1. INTRODUCCIÓN ══════════════════════════════════════════
\section{Introducci\'on y Descripci\'on del Problema}

\subsection{Contexto del Proyecto}

EcoMarket es una plataforma de comercio electr\'onico dise\~nada para la regi\'on de
Puno que conecta a productores locales de productos ecol\'ogicos y org\'anicos con
consumidores conscientes del medio ambiente. La regi\'on de Puno, ubicada en el
altiplano peruano a orillas del lago Titicaca, posee una rica tradici\'on agr\'icola basada
en cultivos nativos como la quinua, ca\~nihua, kiwicha y papas nativas, producidos de
manera tradicional por peque\~nos agricultores.

Sin embargo, estos productores enfrentan serias dificultades para acceder a mercados
m\'as amplios. EcoMarket surge como soluci\'on integral que facilita la venta en l\'inea e
incorpora un sistema de \textbf{auditor\'ia qu\'imica} basado en Inteligencia Artificial que
analiza los ingredientes de cada producto y emite certificados de pureza ecol\'ogica
con puntuaciones \textbf{Eco-Score}, generando confianza tanto en compradores locales
como en mercados internacionales.

\subsection{Problem\'atica Identificada}

\begin{itemize}[leftmargin=1.5cm]
  \item \textbf{Acceso limitado al mercado.} Los productores ecol\'ogicos dependen de
    intermediarios que reducen sus m\'argenes de ganancia. No cuentan con canales
    digitales para llegar directamente al consumidor final.

  \item \textbf{Desconfianza del consumidor.} No existe un mecanismo verificable que
    certifique la autenticidad de los productos org\'anicos. El consumidor no puede
    distinguir entre productos genuinamente ecol\'ogicos y los que solo se anuncian
    como tales.

  \item \textbf{Falta de digitalizaci\'on.} Los peque\~nos productores agr\'icolas no cuentan
    con herramientas tecnol\'ogicas accesibles para gestionar ventas, inventario y
    pagos.

  \item \textbf{M\'etodos de pago inadecuados.} Las plataformas existentes no integran
    m\'etodos como Yape y Plin, ampliamente utilizados en Puno donde la
    bancarizaci\'on es limitada.

  \item \textbf{Ausencia de trazabilidad qu\'imica.} Ning\'un sistema analiza autom\'aticamente
    los ingredientes de los productos agr\'icolas ni certifica su inocuidad, lo que
    representa un riesgo para los consumidores.
\end{itemize}

\subsection{Objetivo General}

Desarrollar una plataforma de comercio electr\'onico para productos ecol\'ogicos de la
regi\'on de Puno que integre un sistema de auditor\'ia qu\'imica mediante IA para
garantizar la autenticidad y calidad de los productos, conectando directamente a
productores locales con consumidores nacionales e internacionales.

\subsection{Objetivos Espec\'ificos}

\begin{enumerate}[leftmargin=1.5cm]
  \item Dise\~nar un cat\'alogo digital de productos ecol\'ogicos con filtros por categor\'ia,
    Eco-Score y ubicaci\'on geogr\'afica.
  \item Desarrollar un m\'odulo de autenticaci\'on seguro con roles diferenciados:
    Consumidor, Productor y Administrador.
  \item Implementar un motor de auditor\'ia qu\'imica que calcule el Eco-Score (0--100)
    de cada producto mediante IA.
  \item Integrar m\'etodos de pago locales (Yape, Plin) e internacionales (Stripe).
  \item Generar certificados PDF de auditor\'ia con hash SHA-256 para garantizar la
    inmutabilidad de los resultados.
  \item Dise\~nar una arquitectura de microservicios escalable y desplegable en la nube.
\end{enumerate}

\newpage

% ══ 2. ENTREVISTAS ═══════════════════════════════════════════
\section{Entrevistas}

Las entrevistas se realizaron a tres actores clave del ecosistema de productos
ecol\'ogicos en la regi\'on de Puno: un productor agr\'icola, un consumidor local y una
comerciante del mercado central. El objetivo fue identificar las necesidades,
expectativas y requerimientos funcionales de la plataforma EcoMarket.

% ── Entrevista 1 ─────────────────────────────────────────────
\subsection{Entrevista N.\textdegree{}1 --- Productor Ecol\'ogico}

\begin{center}
\begin{tabular}{|L{3.8cm}|L{9.7cm}|}
\hline
\rowcolor{ecogreen}
\multicolumn{2}{|c|}{\color{white}\textbf{Ficha de Entrevista}} \\
\hline
\textbf{Entrevistado:}  & Juan Quispe Callata \\
\hline
\textbf{Rol:}           & Agricultor ecol\'ogico, Comunidad de Chucuito, Puno \\
\hline
\textbf{Productos:}     & Quinua real, ca\~nihua, papas nativas \\
\hline
\textbf{Fecha:}         & 5 de marzo de 2026 \\
\hline
\textbf{Modalidad:}     & Presencial en la Comunidad de Chucuito \\
\hline
\textbf{Duraci\'on:}    & 40 minutos \\
\hline
\textbf{Entrevistador:} & Garate Quispe Josep Vladimir \\
\hline
\end{tabular}
\end{center}

\subsubsection{Transcripci\'on Resumida}

\textbf{Entrevistador:} Don Juan, \textquestiondown c\'omo vende actualmente sus productos ecol\'ogicos?

\textbf{Productor:} Vendo en la feria sabatina de Puno y a veces a intermediarios que
llevan mis productos a Juliaca y Arequipa. El problema es que me pagan muy poco.
La quinua real me compran a 3 soles el kilo y ellos la venden a 12 soles en la ciudad.
No tengo forma de llegar directamente al consumidor final.

\textbf{Entrevistador:} \textquestiondown Qu\'e opina de vender sus productos por internet?

\textbf{Productor:} Me gustar\'ia, pero no s\'e c\'omo hacerlo. Mis hijos me ense\~naron a
usar WhatsApp, pero una tienda virtual me parece complicado. Adem\'as, la gente
desconf\'ia, quieren saber si realmente es org\'anico. Yo tengo mi certificaci\'on del
SENASA, pero no s\'e c\'omo mostrarla en internet.

\textbf{Entrevistador:} \textquestiondown Qu\'e tan importante es la certificaci\'on ecol\'ogica para usted?

\textbf{Productor:} Es muy importante. Yo no uso pesticidas ni fertilizantes qu\'imicos,
todo es natural como se ha hecho siempre en mi comunidad. Si la plataforma pudiera
analizar mis productos y dar un puntaje de confianza, ser\'ia bueno para
diferenciarme de los que no son genuinos.

\textbf{Entrevistador:} \textquestiondown Qu\'e funcionalidades necesita en la plataforma?

\textbf{Productor:} Primero, que sea f\'acil de usar desde mi celular y pueda subir fotos
de mis productos. Segundo, que los clientes vean mi certificaci\'on y el an\'alisis.
Tercero, recibir pagos directos sin intermediarios. Y cuarto, que me ayude a
coordinar las entregas.

\subsubsection{Hallazgos Clave}

\begin{itemize}[leftmargin=1.5cm]
  \item Proceso de venta 100\,\% dependiente de intermediarios que reducen m\'argenes
    de ganancia en m\'as del 70\,\%.
  \item Necesidad de interfaz m\'ovil simple e intuitiva accesible desde smartphones
    de gama baja.
  \item Demanda de un sistema de certificaci\'on digital verificable y con respaldo
    de entidades como SENASA.
  \item Preferencia por pagos directos sin comisiones de intermediarios.
  \item Inters en coordinar log\'istica de entrega desde la propia plataforma.
\end{itemize}

% ── Entrevista 2 ─────────────────────────────────────────────
\subsection{Entrevista N.\textdegree{}2 --- Consumidora Local}

\begin{center}
\begin{tabular}{|L{3.8cm}|L{9.7cm}|}
\hline
\rowcolor{ecogreen}
\multicolumn{2}{|c|}{\color{white}\textbf{Ficha de Entrevista}} \\
\hline
\textbf{Entrevistado:}  & Mar\'ia Mamani Huanca \\
\hline
\textbf{Rol:}           & Profesora de nivel primario, consumidora local \\
\hline
\textbf{Lugar:}         & Ciudad de Puno \\
\hline
\textbf{Fecha:}         & 8 de marzo de 2026 \\
\hline
\textbf{Modalidad:}     & Presencial \\
\hline
\textbf{Duraci\'on:}    & 30 minutos \\
\hline
\textbf{Entrevistador:} & Quispe Huanca Arnold Diego \\
\hline
\end{tabular}
\end{center}

\subsubsection{Transcripci\'on Resumida}

\textbf{Entrevistador:} Mar\'ia, \textquestiondown d\'onde compra actualmente sus alimentos?

\textbf{Consumidora:} Compro en el mercado central y en algunas bodegas del barrio.
Me gustar\'ia comprar productos ecol\'ogicos porque he le\'ido que son m\'as saludables,
pero en el mercado no siempre s\'e cu\'al es realmente org\'anico. Todo se ve igual y
a veces los precios son m\'as altos sin saber si vale la pena.

\textbf{Entrevistador:} \textquestiondown Le gustar\'ia una plataforma con certificaci\'on ecol\'ogica?

\textbf{Consumidora:} S\'i, definitivamente. Sobre todo si puedo ver el an\'alisis de cada
producto, su puntuaci\'on ecol\'ogica y qu\'e ingredientes tiene. Tambi\'en me gustar\'ia
saber de d\'onde viene, porque prefiero apoyar a los agricultores de la regi\'on. Si
tiene entregas a domicilio, mucho mejor.

\textbf{Entrevistador:} \textquestiondown Qu\'e m\'etodos de pago prefiere?

\textbf{Consumidora:} Uso Yape y Plin casi a diario. Tambi\'en tengo tarjeta de d\'ebito.
Ser\'ia bueno tener varias opciones. A veces me da desconfianza poner datos de
tarjeta en internet, pero si la plataforma tiene buena reputaci\'on lo usar\'ia.

\textbf{Entrevistador:} \textquestiondown Qu\'e otra funci\'on le gustar\'ia?

\textbf{Consumidora:} Ver el historial de mis compras anteriores y quiz\'as acumular
puntos por cada compra para obtener descuentos. En otras tiendas lo hacen y es
muy \'util.

\subsubsection{Hallazgos Clave}

\begin{itemize}[leftmargin=1.5cm]
  \item Existe demanda insatisfecha de productos ecol\'ogicos verificados en Puno.
  \item Los consumidores valoran conocer la procedencia geogr\'afica del producto.
  \item Yape y Plin son los m\'etodos de pago preferidos en la regi\'on.
  \item El Eco-Score y los certificados de auditor\'ia generan confianza en la compra.
  \item Inters en historial de pedidos y sistema de fidelizaci\'on por puntos.
\end{itemize}

% ── Entrevista 3 ─────────────────────────────────────────────
\subsection{Entrevista N.\textdegree{}3 --- Comerciante del Mercado Central}

\begin{center}
\begin{tabular}{|L{3.8cm}|L{9.7cm}|}
\hline
\rowcolor{ecogreen}
\multicolumn{2}{|c|}{\color{white}\textbf{Ficha de Entrevista}} \\
\hline
\textbf{Entrevistado:}  & Rosa Condori Vilca \\
\hline
\textbf{Rol:}           & Comerciante, Mercado Central de Puno \\
\hline
\textbf{Productos:}     & Quinuas, habas, cebada, quesos regionales \\
\hline
\textbf{Fecha:}         & 10 de marzo de 2026 \\
\hline
\textbf{Modalidad:}     & Presencial en el Mercado Central \\
\hline
\textbf{Duraci\'on:}    & 25 minutos \\
\hline
\textbf{Entrevistador:} & Garate Quispe Josep Vladimir \\
\hline
\end{tabular}
\end{center}

\subsubsection{Transcripci\'on Resumida}

\textbf{Entrevistador:} Sra. Rosa, \textquestiondown ha considerado vender sus productos en l\'inea?

\textbf{Comerciante:} Lo he pensado, pero no tengo tiempo para aprender plataformas
complicadas. Adem\'as, el env\'io es un problema porque vendo productos frescos. Me
preocupa que los clientes no paguen o me estafen. Prefiero el trato directo,
cara a cara.

\textbf{Entrevistador:} \textquestiondown Cree que una plataforma digital aumentar\'ia sus ventas?

\textbf{Comerciante:} Si es f\'acil de usar y tiene buena reputaci\'on, s\'i. Sobre todo si
atrae a turistas o gente de afuera de Puno. Muchos turistas preguntan por quinua
real, ca\~nihua y quesos de cabra, pero no siempre encuentran d\'onde comprar con
confianza.

\textbf{Entrevistador:} \textquestiondown Qu\'e opina de un sistema de auditor\'ia de ingredientes?

\textbf{Comerciante:} Me parece excelente. As\'i el cliente puede ver que el producto es
realmente org\'anico y no solo publicidad. Yo conozco personalmente a los
productores. Si la plataforma puede certificar eso con un an\'alisis, ser\'ia un gran
respaldo para los comerciantes honestos.

\subsubsection{Hallazgos Clave}

\begin{itemize}[leftmargin=1.5cm]
  \item Los comerciantes necesitan interfaces simples que no requieran formaci\'on t\'ecnica.
  \item Existe demanda tur\'istica de productos t\'ipicos regionales con calidad certificada.
  \item La auditor\'ia qu\'imica es vista como un diferenciador competitivo frente a otras
    plataformas.
  \item Se requiere un sistema de pagos seguro que proteja contra fraudes en l\'inea.
\end{itemize}

\newpage

% ══ 3. REQUISITOS ════════════════════════════════════════════
\section{An\'alisis de Requisitos}

\subsection{Requisitos Funcionales}

\begin{center}
\begin{longtable}{|C{1.6cm}|L{4cm}|L{8.4cm}|}
\hline
\rowcolor{ecogreen}
\color{white}\textbf{ID} &
\color{white}\textbf{Nombre} &
\color{white}\textbf{Descripci\'on} \\
\hline
\endfirsthead
\hline
\rowcolor{ecogreen}
\color{white}\textbf{ID} &
\color{white}\textbf{Nombre} &
\color{white}\textbf{Descripci\'on} \\
\hline
\endhead

\rowcolor{ecogreenlight}
\textbf{RF-01} & Registro de Usuario &
El usuario se registra con nombre, email, contrase\~na y tel\'efono. Los
Productores a\~naden RUC y certificaciones ecol\'ogicas. \\
\hline
RF-02 & Inicio de Sesi\'on &
El usuario se autentica con email y contrase\~na. El sistema genera un
token JWT v\'alido por 24 horas. \\
\hline
\rowcolor{ecogreenlight}
\textbf{RF-03} & Gesti\'on de Productos &
El Productor crea, edita y gestiona productos: im\'agenes, ingredientes,
precio, stock y coordenadas GPS del origen. \\
\hline
RF-04 & Cat\'alogo con Filtros &
El visitante navega el cat\'alogo con filtros por categor\'ia, texto y
Eco-Score m\'inimo. Solo se muestran productos aprobados. \\
\hline
\rowcolor{ecogreenlight}
\textbf{RF-05} & Auditor\'ia Qu\'imica IA &
El Administrador ejecuta el an\'alisis de ingredientes. La IA calcula el
Eco-Score (0--100): APROBADO ($\geq$70) o RECHAZADO ($<$70). \\
\hline
RF-06 & Certificado PDF &
El sistema genera un certificado PDF con Eco-Score, tabla de ingredientes,
badges de calidad y hash SHA-256 de verificaci\'on. \\
\hline
\rowcolor{ecogreenlight}
\textbf{RF-07} & Carrito de Compras &
El cliente agrega, elimina y modifica productos. El sistema valida el
stock disponible en tiempo real. \\
\hline
RF-08 & Proceso de Pago &
El cliente completa la compra v\'ia Yape, Plin, TuPay (locales) o Stripe
(internacional). Se genera una orden confirmada. \\
\hline
\rowcolor{ecogreenlight}
\textbf{RF-09} & Gesti\'on de Usuarios &
El Administrador activa/desactiva cuentas, verifica RUC de Productores
y gestiona roles del sistema. \\
\hline
RF-10 & Dashboard Ventas &
El Productor visualiza estad\'isticas propias. El Administrador accede
a m\'etricas globales del sistema. \\
\hline
\rowcolor{ecogreenlight}
\textbf{RF-11} & Notificaciones Push &
El sistema notifica autom\'aticamente sobre pagos confirmados, auditor\'ias
completadas y cambios de estado de pedidos. \\
\hline
RF-12 & Historial de Pedidos &
El cliente consulta su historial de compras con detalle de productos,
importes y estados. \\
\hline
\rowcolor{ecogreenlight}
\textbf{RF-13} & Badges Ecol\'ogicos &
El motor de auditor\'ia asigna badges: Eco-Friendly, Toxic-Free, Vegan,
Cruelty Free y Plastic Free. \\
\hline
RF-14 & Verificaci\'on Hash &
Cualquier usuario puede verificar la autenticidad de un certificado PDF
ingresando su hash SHA-256. \\
\hline
\rowcolor{ecogreenlight}
\textbf{RF-15} & Panel de Administraci\'on &
El Administrador gestiona usuarios, productos, auditor\'ias y configuraciones
del sistema desde un panel centralizado. \\
\hline
\end{longtable}
\end{center}

\subsection{Requisitos No Funcionales}

\begin{center}
\begin{longtable}{|C{1.6cm}|L{3cm}|L{9.4cm}|}
\hline
\rowcolor{ecogreen}
\color{white}\textbf{ID} &
\color{white}\textbf{Categor\'ia} &
\color{white}\textbf{Descripci\'on} \\
\hline
\endfirsthead
\hline
\rowcolor{ecogreen}
\color{white}\textbf{ID} &
\color{white}\textbf{Categor\'ia} &
\color{white}\textbf{Descripci\'on} \\
\hline
\endhead

\rowcolor{ecogreenlight}
\textbf{RNF-01} & Rendimiento &
El frontend carga en menos de 3 segundos en 4G. Las APIs responden
en menos de 500~ms para operaciones CRUD est\'andar. \\
\hline
RNF-02 & Disponibilidad &
El sistema estar\'a disponible el 99\,\% del tiempo. Vercel garantiza
uptime del frontend; Render gestiona los microservicios backend. \\
\hline
\rowcolor{ecogreenlight}
\textbf{RNF-03} & Seguridad &
Autenticaci\'on mediante JWT con expiraci\'on de 24 horas. Contrase\~nas
cifradas con BCrypt. Comunicaciones bajo HTTPS/TLS~1.3. \\
\hline
RNF-04 & Responsividad &
La interfaz ser\'a completamente funcional en dispositivos m\'oviles desde
320~px de ancho, usando Tailwind CSS con dise\~no responsivo. \\
\hline
\rowcolor{ecogreenlight}
\textbf{RNF-05} & Mantenibilidad &
Arquitectura de microservicios con responsabilidades separadas. Cada
servicio es independiente y desplegable aut\'onomamente. \\
\hline
RNF-06 & Escalabilidad &
Cada microservicio puede escalarse horizontalmente de forma independiente
mediante contenedores Docker en Render. \\
\hline
\rowcolor{ecogreenlight}
\textbf{RNF-07} & Usabilidad &
La interfaz del cliente ser\'a intuitiva para personas sin conocimientos
t\'ecnicos. Tiempo de aprendizaje m\'aximo: 5 minutos. \\
\hline
\end{longtable}
\end{center}

\newpage

% ══ 4. CASOS DE USO ══════════════════════════════════════════
\section{Casos de Uso}

\subsection{Actores del Sistema}

\begin{center}
\begin{tabular}{|L{2.8cm}|L{11.2cm}|}
\hline
\rowcolor{ecogreen}
\color{white}\textbf{Actor} & \color{white}\textbf{Descripci\'on} \\
\hline
\textbf{Cliente} &
Usuario final que navega el cat\'alogo, compra productos y visualiza
certificados de auditor\'ia. Puede ser visitante an\'onimo o registrado. \\
\hline
\textbf{Productor} &
Vendedor que registra su cat\'alogo ecol\'ogico, visualiza resultados de
auditor\'ias y estad\'isticas de ventas. \\
\hline
\textbf{Administrador} &
Rol con acceso total. Gestiona usuarios, ejecuta auditor\'ias qu\'imicas,
aprueba/rechaza productos y genera reportes globales. \\
\hline
\textbf{Sistema} &
Agente interno que ejecuta el motor de IA, genera certificados PDF,
env\'ia notificaciones y gestiona cambios de estado autom\'aticamente. \\
\hline
\end{tabular}
\end{center}

\subsection{Relaci\'on Actores --- Casos de Uso}

\begin{center}
\begin{tabular}{|L{4cm}|C{2cm}|L{8cm}|}
\hline
\rowcolor{ecogreen}
\color{white}\textbf{Actor} & \color{white}\textbf{ID} & \color{white}\textbf{Caso de Uso} \\
\hline
Cliente                     & CU-01 & Registro de Usuario \\
\hline
Cliente / Productor / Admin & CU-02 & Inicio de Sesi\'on \\
\hline
Productor                   & CU-03 & Gesti\'on de Cat\'alogo de Productos \\
\hline
Cliente                     & CU-04 & Ver Cat\'alogo y Filtrar Productos \\
\hline
Administrador / Sistema     & CU-05 & Auditor\'ia Qu\'imica de Producto \\
\hline
Admin / Cliente             & CU-06 & Generaci\'on de Certificado PDF \\
\hline
Cliente                     & CU-07 & Compra de Productos \\
\hline
Administrador               & CU-08 & Gesti\'on de Usuarios \\
\hline
Productor / Admin           & CU-09 & Dashboard de Ventas \\
\hline
Sistema                     & CU-10 & Notificaciones Push \\
\hline
\end{tabular}
\end{center}

\subsection{Especificaciones Detalladas de Casos de Uso}

% ── CU-01 ────────────────────────────────────────────────────
\begin{center}
\begin{tabular}{|L{3.5cm}|L{10.5cm}|}
\hline
\rowcolor{ecogreen}
\multicolumn{2}{|c|}{\color{white}\textbf{CU-01 --- Registro de Usuario}} \\
\hline
\textbf{Actores:}        & Cliente, Productor \\
\hline
\textbf{Precondici\'on:} & El email no debe estar registrado previamente. \\
\hline
\textbf{Postcondici\'on:}& Usuario registrado con estado \texttt{activo} y token JWT generado. \\
\hline
\multicolumn{2}{|l|}{\cellcolor{ecogreenlight}\textbf{Flujo Principal}} \\
\hline
\multicolumn{2}{|L{14cm}|}{
1.~El usuario accede al formulario y selecciona su rol (Cliente o Productor).\\
2.~Completa: nombre, email, contrase\~na y tel\'efono.\\
3.~Los Productores a\~naden RUC y certificaciones ecol\'ogicas.\\
4.~El sistema valida los datos y verifica que el email no est\'e registrado.\\
5.~Crea la cuenta y retorna un token JWT para inicio de sesi\'on autom\'atico.
} \\
\hline
\multicolumn{2}{|l|}{\cellcolor{ecogreenlight}\textbf{Flujos Alternativos}} \\
\hline
\multicolumn{2}{|L{14cm}|}{
\textbf{A1:} Si el email ya existe, muestra: ``El email ya est\'a registrado''.\\
\textbf{A2:} Si la contrase\~na no cumple requisitos de seguridad, se muestra un aviso descriptivo.
} \\
\hline
\end{tabular}
\end{center}

\vspace{0.4cm}

% ── CU-02 ────────────────────────────────────────────────────
\begin{center}
\begin{tabular}{|L{3.5cm}|L{10.5cm}|}
\hline
\rowcolor{ecogreen}
\multicolumn{2}{|c|}{\color{white}\textbf{CU-02 --- Inicio de Sesi\'on}} \\
\hline
\textbf{Actores:}        & Cliente, Productor, Administrador \\
\hline
\textbf{Precondici\'on:} & El usuario debe estar registrado y activo. \\
\hline
\textbf{Postcondici\'on:}& Token JWT v\'alido por 24 horas con acceso seg\'un rol. \\
\hline
\multicolumn{2}{|l|}{\cellcolor{ecogreenlight}\textbf{Flujo Principal}} \\
\hline
\multicolumn{2}{|L{14cm}|}{
1.~El usuario ingresa email y contrase\~na.\\
2.~El sistema valida las credenciales contra la base de datos.\\
3.~Si son correctas, genera un token JWT con roles y permisos.\\
4.~Redirige al dashboard correspondiente seg\'un el rol del usuario.
} \\
\hline
\multicolumn{2}{|l|}{\cellcolor{ecogreenlight}\textbf{Flujos Alternativos}} \\
\hline
\multicolumn{2}{|L{14cm}|}{
\textbf{A1:} Credenciales incorrectas: el servidor retorna error 401.\\
\textbf{A2:} Cuenta inactiva: el sistema informa que contacte al administrador.
} \\
\hline
\end{tabular}
\end{center}

\vspace{0.4cm}

% ── CU-03 ────────────────────────────────────────────────────
\begin{center}
\begin{tabular}{|L{3.5cm}|L{10.5cm}|}
\hline
\rowcolor{ecogreen}
\multicolumn{2}{|c|}{\color{white}\textbf{CU-03 --- Gesti\'on de Cat\'alogo de Productos}} \\
\hline
\textbf{Actores:}        & Productor \\
\hline
\textbf{Precondici\'on:} & Productor autenticado con RUC verificado. \\
\hline
\textbf{Postcondici\'on:}& Producto almacenado con estado \texttt{PENDIENTE}. Administrador notificado. \\
\hline
\multicolumn{2}{|l|}{\cellcolor{ecogreenlight}\textbf{Flujo Principal}} \\
\hline
\multicolumn{2}{|L{14cm}|}{
1.~El Productor accede al panel y selecciona ``Nuevo Producto''.\\
2.~Completa: nombre, descripci\'on, categor\'ia, precio, stock e im\'agenes.\\
3.~Registra los ingredientes del producto y el origen GPS.\\
4.~El sistema guarda con estado \texttt{PENDIENTE} y notifica al Administrador.\\
5.~El producto solo es visible para el Productor hasta ser aprobado.
} \\
\hline
\multicolumn{2}{|l|}{\cellcolor{ecogreenlight}\textbf{Flujos Alternativos}} \\
\hline
\multicolumn{2}{|L{14cm}|}{
\textbf{A1:} Si falta un campo obligatorio, el sistema muestra validaci\'on en l\'inea.\\
\textbf{A2:} Si la imagen supera 5~MB, solicita una imagen m\'as peque\~na.
} \\
\hline
\end{tabular}
\end{center}

\vspace{0.4cm}

% ── CU-05 ────────────────────────────────────────────────────
\begin{center}
\begin{tabular}{|L{3.5cm}|L{10.5cm}|}
\hline
\rowcolor{ecogreen}
\multicolumn{2}{|c|}{\color{white}\textbf{CU-05 --- Auditor\'ia Qu\'imica de Producto}} \\
\hline
\textbf{Actores:}        & Administrador, Sistema (IA) \\
\hline
\textbf{Precondici\'on:} & Producto en estado \texttt{PENDIENTE} con ingredientes registrados. \\
\hline
\textbf{Postcondici\'on:}& Producto APROBADO (Eco-Score $\geq$70) o RECHAZADO ($<$70). Certificado PDF con SHA-256 generado. \\
\hline
\multicolumn{2}{|l|}{\cellcolor{ecogreenlight}\textbf{Flujo Principal}} \\
\hline
\multicolumn{2}{|L{14cm}|}{
1.~El Administrador selecciona un producto pendiente de auditor\'ia.\\
2.~El sistema env\'ia los ingredientes al motor de IA para an\'alisis.\\
3.~La IA eval\'ua cada ingrediente contra la BD de sustancias prohibidas.\\
4.~Calcula el Eco-Score restando puntos por sustancias nocivas.\\
5.~Asigna badges: Eco-Friendly, Toxic-Free, Vegan, Cruelty Free, Plastic Free.\\
6.~Eco-Score $\geq$ 70: estado \texttt{APROBADO}; $<$ 70: estado \texttt{RECHAZADO}.\\
7.~Genera certificado PDF con hash SHA-256 y notifica al Productor.
} \\
\hline
\multicolumn{2}{|l|}{\cellcolor{ecogreenlight}\textbf{Flujos Alternativos}} \\
\hline
\multicolumn{2}{|L{14cm}|}{
\textbf{A1:} Sin ingredientes registrados: el sistema solicita completar la informaci\'on.\\
\textbf{A2:} Motor de IA no disponible: el sistema programa reintento autom\'atico.
} \\
\hline
\end{tabular}
\end{center}

\vspace{0.4cm}

% ── CU-07 ────────────────────────────────────────────────────
\begin{center}
\begin{tabular}{|L{3.5cm}|L{10.5cm}|}
\hline
\rowcolor{ecogreen}
\multicolumn{2}{|c|}{\color{white}\textbf{CU-07 --- Compra de Productos}} \\
\hline
\textbf{Actores:}        & Cliente \\
\hline
\textbf{Precondici\'on:} & Cliente autenticado. Productos aprobados con stock disponible. \\
\hline
\textbf{Postcondici\'on:}& Orden con estado \texttt{CONFIRMADO}. Stock descontado. Notificaciones enviadas. \\
\hline
\multicolumn{2}{|l|}{\cellcolor{ecogreenlight}\textbf{Flujo Principal}} \\
\hline
\multicolumn{2}{|L{14cm}|}{
1.~El cliente navega el cat\'alogo y agrega productos al carrito.\\
2.~Procede al checkout y revisa el resumen del pedido.\\
3.~Selecciona el m\'etodo de pago: Yape, Plin, TuPay o Stripe.\\
4a. Para m\'etodos locales: el sistema genera c\'odigo de referencia y confirma.\\
4b. Para Stripe: redirige a la pasarela segura y espera el webhook de confirmaci\'on.\\
5.~El sistema crea la orden, descuenta el stock y env\'ia notificaciones.
} \\
\hline
\multicolumn{2}{|l|}{\cellcolor{ecogreenlight}\textbf{Flujos Alternativos}} \\
\hline
\multicolumn{2}{|L{14cm}|}{
\textbf{A1:} Si el pago falla, el sistema informa al cliente y no crea la orden.\\
\textbf{A2:} Si el stock se agota entre selecci\'on y pago, el sistema notifica al cliente.
} \\
\hline
\end{tabular}
\end{center}

\vspace{0.4cm}

% ── CU-08 ────────────────────────────────────────────────────
\begin{center}
\begin{tabular}{|L{3.5cm}|L{10.5cm}|}
\hline
\rowcolor{ecogreen}
\multicolumn{2}{|c|}{\color{white}\textbf{CU-08 --- Gesti\'on de Usuarios}} \\
\hline
\textbf{Actores:}        & Administrador \\
\hline
\textbf{Precondici\'on:} & Administrador autenticado con rol \texttt{ADMIN}. \\
\hline
\textbf{Postcondici\'on:}& Usuario actualizado. Log de auditor\'ia registrado. \\
\hline
\multicolumn{2}{|l|}{\cellcolor{ecogreenlight}\textbf{Flujo Principal}} \\
\hline
\multicolumn{2}{|L{14cm}|}{
1.~El Administrador accede al panel y visualiza la lista de usuarios.\\
2.~Filtra por rol (CLIENTE, PRODUCTOR) y estado (activo/inactivo).\\
3.~Selecciona un usuario para ver su perfil completo.\\
4.~Puede: activar/desactivar cuenta, verificar RUC o asignar nuevo rol.\\
5.~El sistema aplica el cambio y registra la acci\'on en el log de auditor\'ia.
} \\
\hline
\multicolumn{2}{|l|}{\cellcolor{ecogreenlight}\textbf{Flujos Alternativos}} \\
\hline
\multicolumn{2}{|L{14cm}|}{
\textbf{A1:} Si el usuario no tiene rol ADMIN, el servidor retorna error 403.
} \\
\hline
\end{tabular}
\end{center}

\newpage

% ══ 5. ARQUITECTURA ══════════════════════════════════════════
\section{Arquitectura del Sistema}

\subsection{Patr\'on Arquitect\'onico: Microservicios}

EcoMarket implementa una arquitectura de \textbf{microservicios}, donde cada funci\'on opera
como un servicio independiente con su propia base de datos y tecnolog\'ia. Un
\textbf{API Gateway} centraliza las solicitudes del frontend y las enruta al servicio correspondiente, gestionando CORS de forma centralizada.

\medskip
Flujo: \texttt{Frontend (Vercel)} $\longrightarrow$ \texttt{API Gateway :8080}
$\longrightarrow$ \texttt{Microservicio espec\'ifico} $\longrightarrow$ \texttt{Base de datos}

\begin{center}
\begin{tabular}{|L{3.2cm}|L{4.5cm}|L{5.8cm}|}
\hline
\rowcolor{ecogreen}
\color{white}\textbf{Capa} &
\color{white}\textbf{Descripci\'on} &
\color{white}\textbf{Tecnolog\'ias} \\
\hline
\textbf{Presentaci\'on} &
Interfaz de usuario, componentes React y gesti\'on del carrito. &
Next.js 16, React 19, Tailwind CSS, Zustand \\
\hline
\textbf{Gateway} &
Enrutamiento de solicitudes y CORS centralizado. &
Node.js HTTP nativo, puerto 8080 \\
\hline
\textbf{Servicios Core} &
Autenticaci\'on, autorizaci\'on y cat\'alogo de productos. &
Java 17, Spring Boot 3.2, Spring Security \\
\hline
\textbf{Servicios Aux.} &
Pagos, auditor\'ia, IA y notificaciones. &
Node.js (Express), Python 3.11 (FastAPI) \\
\hline
\textbf{Persistencia} &
Almacenamiento relacional y documental. &
PostgreSQL (Neon), MongoDB Atlas \\
\hline
\textbf{Infraestructura} &
Contenedorizaci\'on y despliegue en la nube. &
Docker, Render, Vercel \\
\hline
\end{tabular}
\end{center}

\subsection{Stack Tecnol\'ogico}

\begin{center}
\begin{longtable}{|L{3cm}|L{4cm}|L{7cm}|}
\hline
\rowcolor{ecogreen}
\color{white}\textbf{Componente} &
\color{white}\textbf{Tecnolog\'ia} &
\color{white}\textbf{Rol en el Sistema} \\
\hline
\endfirsthead
\hline
\rowcolor{ecogreen}
\color{white}\textbf{Componente} &
\color{white}\textbf{Tecnolog\'ia} &
\color{white}\textbf{Rol en el Sistema} \\
\hline
\endhead

\rowcolor{ecogreenlight}
Frontend        & Next.js 16 (App Router) & SSR/CSR h\'ibrido, enrutamiento y componentes de UI. \\
\hline
UI Framework    & React 19 + Tailwind CSS & Componentes reactivos y dise\~no responsivo. \\
\hline
\rowcolor{ecogreenlight}
Estado Global   & Zustand                 & Carrito de compras con persistencia en localStorage. \\
\hline
API Gateway     & Node.js HTTP            & Enrutamiento a microservicios y CORS centralizado. \\
\hline
\rowcolor{ecogreenlight}
Auth Service    & Java 17 + Spring Boot 3.2 & Registro, login, JWT y control de roles. \\
\hline
Product Service & Java 17 + Spring Boot 3.2 & Cat\'alogo, inventario, \'ordenes y CRUD. \\
\hline
\rowcolor{ecogreenlight}
Payment Service & Node.js + Express       & Integraci\'on Stripe, Yape, Plin y TuPay. \\
\hline
Audit Service   & Python 3.11 + FastAPI   & Auditor\'ia qu\'imica, PDF con ReportLab. \\
\hline
\rowcolor{ecogreenlight}
AI Engine       & Python 3.11 + FastAPI   & An\'alisis de ingredientes y c\'alculo de Eco-Score. \\
\hline
Notif. Service  & Python 3.11 + FastAPI   & Notificaciones push ante eventos del sistema. \\
\hline
\rowcolor{ecogreenlight}
BD Relacional   & PostgreSQL (Neon)       & Usuarios, productos, \'ordenes y auditor\'ias. \\
\hline
BD Documental   & MongoDB Atlas (M0)      & Registros de auditor\'ia y certificados ecol\'ogicos. \\
\hline
\rowcolor{ecogreenlight}
Contenerizaci\'on & Docker                & Cada microservicio empaquetado en imagen Docker. \\
\hline
Deploy Backend  & Render                  & Hosting gratuito con CI/CD autom\'atico desde GitHub. \\
\hline
\rowcolor{ecogreenlight}
Deploy Frontend & Vercel                  & Hosting edge global con CI/CD autom\'atico. \\
\hline
Control Versiones & Git + GitHub          & Repo: \texttt{JosepVl123463/EcoMarketFinalSoftware} \\
\hline
\end{longtable}
\end{center}

\subsection{Control de Acceso Basado en Roles (RBAC)}

El sistema implementa tres roles gestionados mediante JWT generados por el Auth Service.
El API Gateway valida el token en cada solicitud antes de enrutarla.

\begin{center}
\begin{tabular}{|L{2.8cm}|L{11.2cm}|}
\hline
\rowcolor{ecogreen}
\color{white}\textbf{Rol} & \color{white}\textbf{Permisos del Sistema} \\
\hline
\textbf{CLIENTE} &
Ver cat\'alogo, agregar al carrito, comprar productos, ver historial de
pedidos y descargar certificados de auditor\'ia. \\
\hline
\textbf{PRODUCTOR} &
Hereda CLIENTE + gestionar cat\'alogo propio, ver auditor\'ias de sus
productos y acceder al dashboard de ventas. \\
\hline
\textbf{ADMIN} &
Hereda PRODUCTOR + gestionar todos los usuarios, ejecutar auditor\'ias
qu\'imicas, aprobar/rechazar productos y ver reportes globales. \\
\hline
\end{tabular}
\end{center}

\medskip
\textbf{Acceso al sistema en producci\'on:}

\begin{itemize}[leftmargin=1.5cm]
  \item \textbf{URL:} \href{https://eco-market-final-software-s2w6.vercel.app}{\texttt{https://eco-market-final-software-s2w6.vercel.app}}
  \item \textbf{Email administrador:} \texttt{admin@ecomarket.pe}
  \item \textbf{Contrase\~na:} \texttt{123456789}
\end{itemize}

\newpage

% ══ 6. AVANCE ════════════════════════════════════════════════
\section{Avance de Implementaci\'on}

El sistema tiene el flujo principal completamente implementado y desplegado en
producci\'on, cubriendo el ciclo completo: registro, cat\'alogo, auditor\'ia qu\'imica,
compra y notificaciones.

\begin{center}
\begin{longtable}{|L{7.5cm}|C{2.2cm}|L{4.3cm}|}
\hline
\rowcolor{ecogreen}
\color{white}\textbf{Funcionalidad} &
\color{white}\textbf{Estado} &
\color{white}\textbf{Servicio / Ruta} \\
\hline
\endfirsthead
\hline
\rowcolor{ecogreen}
\color{white}\textbf{Funcionalidad} &
\color{white}\textbf{Estado} &
\color{white}\textbf{Servicio / Ruta} \\
\hline
\endhead

Registro e inicio de sesi\'on con JWT              & \textbf{Completo}   & auth-service \\
\hline
\rowcolor{ecogreenlight}
Roles diferenciados (CLIENTE, PRODUCTOR, ADMIN)    & \textbf{Completo}   & auth-service \\
\hline
Cat\'alogo de productos con filtros                & \textbf{Completo}   & product-service \\
\hline
\rowcolor{ecogreenlight}
Gesti\'on de productos por el Productor            & \textbf{Completo}   & product-service \\
\hline
Motor de auditor\'ia qu\'imica con IA              & \textbf{Completo}   & ai-engine \\
\hline
\rowcolor{ecogreenlight}
C\'alculo de Eco-Score (0--100)                    & \textbf{Completo}   & ai-engine \\
\hline
Asignaci\'on de badges ecol\'ogicos                & \textbf{Completo}   & ai-engine \\
\hline
\rowcolor{ecogreenlight}
Certificado PDF con hash SHA-256                   & \textbf{Completo}   & audit-service \\
\hline
Registro de auditor\'ias en MongoDB                & \textbf{Completo}   & audit-service \\
\hline
\rowcolor{ecogreenlight}
Carrito de compras con Zustand                     & \textbf{Completo}   & Frontend \\
\hline
Proceso de pago con Stripe                         & \textbf{Completo}   & payment-service \\
\hline
\rowcolor{ecogreenlight}
M\'etodos locales (Yape, Plin, TuPay)              & \textbf{Completo}   & payment-service \\
\hline
API Gateway con enrutamiento y CORS                & \textbf{Completo}   & cloud-gateway \\
\hline
\rowcolor{ecogreenlight}
Panel de administraci\'on                          & \textbf{Completo}   & Frontend /admin \\
\hline
Notificaciones push autom\'aticas                  & \textbf{Completo}   & notification-service \\
\hline
\rowcolor{ecogreenlight}
Despliegue frontend en Vercel                      & \textbf{Completo}   & Vercel CI/CD \\
\hline
Despliegue backend en Render (Docker)              & \textbf{Completo}   & Render + Docker \\
\hline
\rowcolor{ecogreenlight}
Base de datos PostgreSQL en Neon                   & \textbf{Completo}   & Neon serverless \\
\hline
Base de datos MongoDB Atlas                        & \textbf{Completo}   & MongoDB M0 \\
\hline
\rowcolor{lightgray}
Historial de pedidos del cliente                   & \textit{Pendiente}  & product-service \\
\hline
\rowcolor{lightgray}
Sistema de puntos de fidelizaci\'on                & \textit{Pendiente}  & product-service \\
\hline
\rowcolor{lightgray}
Dashboard de m\'etricas avanzadas                  & \textit{Pendiente}  & Frontend \\
\hline
\rowcolor{lightgray}
Mapa de origen del productor (GPS)                 & \textit{Pendiente}  & Frontend \\
\hline
\end{longtable}
\end{center}

% ══ CONCLUSIONES ═════════════════════════════════════════════
\section*{Conclusiones}
\addcontentsline{toc}{section}{Conclusiones}

\begin{itemize}[leftmargin=1.5cm]
  \item Se identificaron y documentaron los principales problemas del ecosistema de
    productores ecol\'ogicos de Puno mediante tres entrevistas estructuradas,
    validando la necesidad de una plataforma digital de comercio ecol\'ogico.

  \item Se definieron 15 requisitos funcionales y 7 no funcionales que cubren el flujo
    completo: desde el registro de productores hasta la compra verificada con
    auditor\'ia qu\'imica mediante IA.

  \item Se dise\~naron 10 casos de uso que describen con precisi\'on las interacciones de
    los actores con el sistema, incluyendo flujos alternativos y condiciones de error.

  \item La arquitectura de microservicios adoptada garantiza separaci\'on de
    responsabilidades, escalabilidad independiente por servicio y facilidad de
    mantenimiento, con un stack moderno desplegado en Vercel y Render.

  \item El avance actual cubre el flujo principal completamente implementado y
    desplegado en producci\'on, incluyendo autenticaci\'on, cat\'alogo, auditor\'ia
    qu\'imica con IA, pagos integrados y notificaciones autom\'aticas.
\end{itemize}

% ══ REFERENCIAS ══════════════════════════════════════════════
\section*{Referencias}
\addcontentsline{toc}{section}{Referencias}

\begin{enumerate}[leftmargin=1.5cm]
  \item Vercel Inc.\ (2024). \textit{Next.js 16 Documentation}.
    Recuperado de \url{https://nextjs.org/docs}

  \item Spring Framework.\ (2024). \textit{Spring Boot 3.2 Reference Documentation}.
    Recuperado de \url{https://docs.spring.io/spring-boot/docs/3.2.x/reference/html/}

  \item FastAPI.\ (2024). \textit{FastAPI Documentation}.
    Recuperado de \url{https://fastapi.tiangolo.com}

  \item Render Inc.\ (2024). \textit{Render Docs --- Docker Deployments}.
    Recuperado de \url{https://render.com/docs}

  \item Neon Inc.\ (2024). \textit{Neon Serverless PostgreSQL Documentation}.
    Recuperado de \url{https://neon.tech/docs}

  \item Sommerville, I.\ (2016). \textit{Software Engineering} (10.\textordfeminine\ ed.).
    Pearson Education.

  \item Pressman, R., \& Maxim, B.\ (2020). \textit{Software Engineering:
    A Practitioner's Approach} (9.\textordfeminine\ ed.). McGraw-Hill Education.

  \item Newman, S.\ (2021). \textit{Building Microservices: Designing Fine-Grained
    Systems} (2.\textordfeminine\ ed.). O'Reilly Media.
\end{enumerate}

\end{document}
